\documentclass[11pt]{article}

\usepackage{amsmath,amssymb,amsthm,mathtools}
\usepackage[margin=1in]{geometry}
\usepackage{microtype}
\usepackage[hidelinks]{hyperref}

\newtheorem{theorem}{Theorem}[section]
\newtheorem{proposition}[theorem]{Proposition}
\newtheorem{lemma}[theorem]{Lemma}
\newtheorem{corollary}[theorem]{Corollary}
\newtheorem{remark}[theorem]{Remark}

\DeclareMathOperator{\Hom}{Hom}
\newcommand{\dd}{\,\mathrm d}

\title{The $(3,5)$ Odd-Walk Inequality\\
Entropy on Finite Graphs, and Relative Entropy on Graphons}
\author{}
\date{}

\begin{document}

\maketitle

\begin{abstract}
Section~\ref{sec:finite} is the Blekherman--Raymond entropy proof of
$t(P_5,W)^3\geq t(P_3,W)^5$ for finite host graphs.  Section~\ref{sec:graphon}
transcribes the same argument --- the same three folds of $P_5$ onto $P_3$, the
same edge multiplicities $(5,5,5)$ and internal-vertex multiplicities
$(1,5,5,1)$ --- directly into the graphon setting, with Shannon entropy against
counting measure replaced by relative entropy against $\mu$.  The transcription
needs no finite host, no sampling and no approximation step, and it is valid on
an arbitrary probability space.  That is what the catalogue requires, and it is
what the Lean formalization follows.
\end{abstract}

\section{The finite case}
\label{sec:finite}

\begin{proposition}[Blekherman--Raymond, $P_3$--$P_5$ case]
Let $P_k$ denote the path with $k$ edges. Then, for every graphon $W$,
\[
    t(P_5,W)^3 \ge t(P_3,W)^5.
\]
\end{proposition}

\begin{proof}
We first prove the inequality for a finite graph $G$ on $n$ vertices.
Write
\[
    \hom(P_k,G):=|\Hom(P_k,G)|.
\]
The desired inequality
\[
    t(P_5,G)^3\ge t(P_3,G)^5
\]
is equivalent to
\[
    n^2 \hom(P_5,G)^3\ge \hom(P_3,G)^5,
\]
since $P_5$ has $6$ vertices and $P_3$ has $4$ vertices.

If $G$ has no edges, then both sides vanish, so there is nothing to prove.
Assume henceforth that $G$ has at least one edge.

Let
\[
    X=(X_0,X_1,X_2,X_3)
\]
be a uniformly random homomorphism from $P_3$ to $G$.
Thus $X$ is a uniformly random walk of length $3$ in $G$, and therefore
\[
    H(X_0,X_1,X_2,X_3)
    =
    \log |\Hom(P_3,G)|.
\]
The base of the logarithm is irrelevant.

We shall construct a random homomorphism from
\[
    F:=P_0\sqcup P_0\sqcup P_5\sqcup P_5\sqcup P_5
\]
to $G$ whose entropy is exactly $5H(X)$.
Since
\[
    |\Hom(F,G)|
    =
    n^2\hom(P_5,G)^3,
\]
the entropy bound
\[
    H(Y)\le \log|\Hom(F,G)|
\]
will then give the result.

\medskip

\noindent
\textbf{Step 1: the Markov property of the uniform random walk.}

The law of $X$ is uniform on all quadruples
$(x_0,x_1,x_2,x_3)$ satisfying
\[
    x_0x_1,\ x_1x_2,\ x_2x_3\in E(G).
\]
Consequently, conditioned on $X_i$, the variables on the two sides of
$X_i$ are conditionally independent.  In particular,
\[
    X_0 \perp (X_2,X_3)\mid X_1,
    \qquad
    (X_0,X_1)\perp X_3\mid X_2.
\]
Hence the entropy of $X$ satisfies the tree entropy identity
\begin{equation}\label{eq:tree-entropy-P3}
\begin{aligned}
    H(X)
    &=
    H(X_0,X_1)
    +H(X_1,X_2)
    +H(X_2,X_3)  \\
    &\qquad
    -H(X_1)-H(X_2).
\end{aligned}
\end{equation}

Indeed, by the chain rule and the conditional independences,
\[
\begin{aligned}
H(X_0,X_1,X_2,X_3)
&=
H(X_1)
+H(X_0\mid X_1)
+H(X_2\mid X_1)
+H(X_3\mid X_2) \\
&=
H(X_0,X_1)+H(X_1,X_2)+H(X_2,X_3)
-H(X_1)-H(X_2).
\end{aligned}
\]

\medskip

\noindent
\textbf{Step 2: three folds of $P_5$ onto $P_3$.}

Label the vertices of $P_3$ by
\[
    0,1,2,3
\]
in their natural order.

Consider the following three graph homomorphisms
$\phi_i:P_5\to P_3$, specified by the images of the six consecutive
vertices of $P_5$:
\[
\begin{aligned}
    \phi_1 &: (0,1,0,1,2,3),\\
    \phi_2 &: (0,1,2,1,2,3),\\
    \phi_3 &: (0,1,2,3,2,3).
\end{aligned}
\]
Thus $\phi_i$ traverses every edge of $P_3$ once, except for the
$i$th edge, which it traverses three times.

We now explain how to ``pull back'' the distribution of $X$ along
each $\phi_i$.

More generally, suppose
\[
    a_0,a_1,\dots,a_5\in\{0,1,2,3\}
\]
form a walk in $P_3$.  Define a random walk
\[
    Y=(Y_0,\dots,Y_5)
\]
in $G$ as follows.  First sample $Y_0$ with the same distribution as
$X_{a_0}$.  Having sampled $Y_j$, sample $Y_{j+1}$ according to the
conditional distribution
\[
    \mathbb P(Y_{j+1}=v\mid Y_j=u)
    :=
    \mathbb P(X_{a_{j+1}}=v\mid X_{a_j}=u).
\]
Then
\[
    Y_j\stackrel{d}{=}X_{a_j}
\]
for every $j$, and
\[
    (Y_j,Y_{j+1})
    \stackrel{d}{=}
    (X_{a_j},X_{a_{j+1}})
\]
for every $j$.
Moreover, $Y$ is supported on $\Hom(P_5,G)$.

Applying this construction to the three walks $\phi_1,\phi_2,\phi_3$
gives random homomorphisms
\[
    Y^{(1)},Y^{(2)},Y^{(3)}\in\Hom(P_5,G).
\]
We take these three random homomorphisms independently.

Finally, let $Z_0$ and $Z_3$ be independent random vertices having,
respectively, the distributions of $X_0$ and $X_3$, and independent
of all the $Y^{(i)}$.

Altogether,
\[
    Y:=(Z_0,Z_3,Y^{(1)},Y^{(2)},Y^{(3)})
\]
is a random element of $\Hom(F,G)$.

\medskip

\noindent
\textbf{Step 3: compute the entropy.}

For any Markov distribution on a path
$v_0v_1\cdots v_m$, one has
\begin{equation}\label{eq:tree-entropy}
    H(V_0,\dots,V_m)
    =
    \sum_{j=0}^{m-1}H(V_j,V_{j+1})
    -
    \sum_{j=1}^{m-1}H(V_j).
\end{equation}
This follows immediately from the chain rule:
\[
\begin{aligned}
H(V_0,\dots,V_m)
&=
H(V_0)+\sum_{j=0}^{m-1}H(V_{j+1}\mid V_j)\\
&=
H(V_0)+\sum_{j=0}^{m-1}
    \bigl(H(V_j,V_{j+1})-H(V_j)\bigr).
\end{aligned}
\]

We apply \eqref{eq:tree-entropy} to the three copies of $P_5$.

Across the three maps $\phi_1,\phi_2,\phi_3$, each edge
\[
    01,\qquad 12,\qquad 23
\]
of $P_3$ occurs exactly five times:
each occurs once in each $\phi_i$, and twice more in the unique
$\phi_i$ which backtracks over that edge.

Hence the positive edge contribution to
\[
    H(Y^{(1)})+H(Y^{(2)})+H(Y^{(3)})
\]
is
\[
    5H(X_0,X_1)
    +5H(X_1,X_2)
    +5H(X_2,X_3).
\]

It remains to count the internal vertices of the three copies of
$P_5$.  Their images under the three $\phi_i$ are
\[
\begin{array}{c|c}
\phi_1 & 1,0,1,2\\
\phi_2 & 1,2,1,2\\
\phi_3 & 1,2,3,2.
\end{array}
\]
Thus, among all internal vertices,
\[
    0\text{ occurs once},\qquad
    1\text{ occurs five times},\qquad
    2\text{ occurs five times},\qquad
    3\text{ occurs once}.
\]
Therefore
\begin{align*}
\sum_{i=1}^3 H(Y^{(i)})
&=
5H(X_0,X_1)+5H(X_1,X_2)+5H(X_2,X_3)\\
&\qquad
-H(X_0)-5H(X_1)-5H(X_2)-H(X_3).
\end{align*}

The two isolated vertices contribute
\[
    H(Z_0)+H(Z_3)=H(X_0)+H(X_3).
\]
Since all five components were chosen independently,
\[
\begin{aligned}
H(Y)
&=
H(Z_0)+H(Z_3)+\sum_{i=1}^3H(Y^{(i)})\\
&=
5\Bigl(
H(X_0,X_1)+H(X_1,X_2)+H(X_2,X_3)
-H(X_1)-H(X_2)
\Bigr).
\end{aligned}
\]
By \eqref{eq:tree-entropy-P3},
\[
    H(Y)=5H(X).
\]
Since $X$ is uniform on $\Hom(P_3,G)$,
\[
    H(Y)
    =
    5\log\hom(P_3,G).
\]

On the other hand, $Y$ is supported on $\Hom(F,G)$, so the
maximum-entropy bound gives
\[
\begin{aligned}
    5\log\hom(P_3,G)
    &=H(Y)\\
    &\le \log|\Hom(F,G)|\\
    &=\log\bigl(n^2\hom(P_5,G)^3\bigr).
\end{aligned}
\]
Exponentiating,
\[
    \hom(P_3,G)^5
    \le
    n^2\hom(P_5,G)^3.
\]
Dividing by $n^{20}$ yields
\[
    t(P_3,G)^5
    \le
    t(P_5,G)^3.
\]

Finally, the inequality extends from finite graphs to graphons by the
standard approximation of graphons by finite graphs together with
continuity of homomorphism densities.  Hence, for every graphon $W$,
\[
    \boxed{t(P_5,W)^3\ge t(P_3,W)^5}.
\]
\end{proof}
\section{The same argument on graphons}
\label{sec:graphon}

The proof of Section~\ref{sec:finite} ends by transferring the finite
inequality to graphons ``by the standard approximation of graphons by finite
graphs together with continuity of homomorphism densities''.  That step is
delicate at the generality the catalogue needs: its graphons live on an
arbitrary probability space, not necessarily atomless or standard Borel, and
the usual step-function approximation is not available there.  The honest
finite-to-graphon route is the $W$-random sampling argument of
\texttt{notes/triangle\_three\_edge\_tail.tex}, which needs a product space, a
conditional Bernoulli construction, a second-moment estimate and convergence in
probability.

None of that is necessary.  The entropy proof transcribes directly.  What
disappears is exactly the appeal to a uniform distribution on a finite set:
Shannon entropy is replaced by relative entropy with respect to $\mu$, the
maximum-entropy bound $H(Y)\leq\log|\Hom(F,G)|$ becomes Gibbs' inequality, and
the two isolated vertices --- present in Section~\ref{sec:finite} only to
supply the factor $n^2$ --- become two applications of the nonnegativity of
relative entropy.

\subsection{Setting}

Let $(\Omega,\mu)$ be an arbitrary probability space and $W\colon\Omega^2\to[0,1]$
a symmetric measurable graphon.  Write $T_Wh(x)=\int W(x,y)h(y)\dd\mu(y)$ and
\[
 d=T_W\mathbf1,\qquad A=T_W^2\mathbf1,\qquad B=T_W^3\mathbf1,
\]
\[
 a_3=t(P_3,W)=\langle\mathbf1,T_W^3\mathbf1\rangle,\qquad
 a_5=t(P_5,W)=\langle\mathbf1,T_W^5\mathbf1\rangle .
\]
Assume until Section~\ref{sec:limit} that $W\geq\varepsilon$ pointwise for some
$\varepsilon>0$; then $d,A,B\geq\varepsilon$, every density below is bounded
above and below by positive constants, and every logarithm appearing is
bounded.  All integrability hypotheses are therefore automatic, and we do not
comment on them again.

\subsection{The $P_3$ measure and its marginals}

Let
\[
 f(x_0,x_1,x_2,x_3)
 =\frac{W(x_0,x_1)W(x_1,x_2)W(x_2,x_3)}{a_3},
\]
a probability density on $\Omega^4$ with respect to $\mu^4$.  This is the
graphon counterpart of ``a uniformly random homomorphism $P_3\to G$''.
Integrating out coordinates gives its one- and two-dimensional marginal
densities in closed form:
\begin{equation}
\begin{gathered}
 m_0=m_3=\frac{B}{a_3},\qquad m_1=m_2=\frac{d\,A}{a_3},\\[2pt]
 M_{01}(u,v)=\frac{W(u,v)A(v)}{a_3},\qquad
 M_{12}(u,v)=\frac{W(u,v)d(u)d(v)}{a_3},\qquad
 M_{23}(u,v)=\frac{W(u,v)A(u)}{a_3}.
\end{gathered}
\label{eq:marginals}
\end{equation}
Each is a probability density: for instance
$\iint M_{01}=a_3^{-1}\int d(v)A(v)\dd\mu(v)
=a_3^{-1}\langle T_W\mathbf1,T_W^2\mathbf1\rangle=1$.
For a reversed orientation write $M_{10}(u,v)=M_{01}(v,u)$, and similarly for
the other edges.

The Markov property of Step~1 of Section~\ref{sec:finite} becomes the exact
factorization
\begin{equation}
 f=\frac{M_{01}M_{12}M_{23}}{m_1\,m_2},
 \label{eq:tree-factorization}
\end{equation}
which one checks by substituting \eqref{eq:marginals} and cancelling
$d(x_1)A(x_1)d(x_2)A(x_2)$.

\subsection{The one analytic input}

\begin{lemma}[Gibbs]
\label{lem:gibbs}
Let $\nu$ be a measure and let $f,g>0$ be probability densities with respect to
$\nu$.  Then $\int f\log(g/f)\dd\nu\leq0$.
\end{lemma}

\begin{proof}
Pointwise $\log u\leq u-1$, so
$f\log(g/f)\leq f\,(g/f-1)=g-f$, and $\int(g-f)\dd\nu=0$.
\end{proof}

Taking $g\equiv1$ gives the only other fact we need:
\begin{equation}
 \int m\log m\dd\nu\ \geq\ 0
 \qquad\text{for every probability density $m$.}
 \label{eq:rel-entropy-nonneg}
\end{equation}
This is what replaces ``$H(Z_0)\leq\log n$'' in Section~\ref{sec:finite}.

\subsection{The three folds}

Let $a=(a_0,\dots,a_5)$ be a walk in $P_3$ and put
\begin{equation}
 q_a(y_0,\dots,y_5)
 =\frac{\prod_{j=0}^{4}M_{a_ja_{j+1}}(y_j,y_{j+1})}
        {\prod_{j=1}^{4}m_{a_j}(y_j)} .
 \label{eq:fold-density}
\end{equation}

\begin{lemma}
\label{lem:fold}
$q_a$ is a probability density on $\Omega^6$; its marginal on the coordinate
pair $(j,j+1)$ is $M_{a_ja_{j+1}}$, and its marginal on the coordinate $j$ is
$m_{a_j}$.
\end{lemma}

\begin{proof}
Integrate out $y_5$ first: $\int M_{a_4a_5}(y_4,y_5)\dd\mu(y_5)=m_{a_4}(y_4)$,
which cancels the factor $m_{a_4}(y_4)$ in the denominator and leaves the same
expression for the walk $(a_0,\dots,a_4)$.  Iterating down to a single
coordinate leaves $m_{a_0}$, of total mass one.  The same computation performed
from either end gives the stated marginals.
\end{proof}

This is \eqref{eq:fold-density} written out; no conditional-probability
machinery is used, which is what makes the argument valid on an arbitrary
probability space.  Set
\[
 E_e=\iint M_e\,\bigl(\log W-\log M_e\bigr)\dd\mu^2
 \quad(e\in\{01,12,23\}),
 \qquad
 V_k=\int m_k\log m_k\dd\mu .
\]
Note $E_e$ does not depend on the orientation of $e$, because $W$ is symmetric.

\begin{lemma}[Tree-entropy identity]
\label{lem:identity}
$E_{01}+E_{12}+E_{23}+V_1+V_2=\log a_3 .$
\end{lemma}

\begin{proof}
Compute $\int f\log f$ in two ways.  From $f=W_{01}W_{12}W_{23}/a_3$ and
Lemma~\ref{lem:fold} applied to the walk $(0,1,2,3)$,
\[
 \int f\log f
 =\sum_e\iint M_e\log W-\log a_3 .
\]
From the factorization \eqref{eq:tree-factorization},
\[
 \int f\log f=\sum_e\iint M_e\log M_e-V_1-V_2 .
\]
Subtract.
\end{proof}

\begin{lemma}[Per-fold bound]
\label{lem:perfold}
For every walk $a$ in $P_3$ of length five,
\[
 \log a_5\ \geq\ \sum_{j=0}^{4}E_{\{a_j,a_{j+1}\}}+\sum_{j=1}^{4}V_{a_j}.
\]
\end{lemma}

\begin{proof}
Apply Lemma~\ref{lem:gibbs} on $\Omega^6$ with $f=q_a$ and
$g=a_5^{-1}\prod_{j=0}^4W(y_j,y_{j+1})$, both probability densities.  Expanding
$\log(g/q_a)$ by \eqref{eq:fold-density} and integrating term by term with
Lemma~\ref{lem:fold},
\[
 0\ \geq\ \int q_a\log\frac{g}{q_a}
 =\sum_{j=0}^{4}\iint M_{a_ja_{j+1}}\bigl(\log W-\log M_{a_ja_{j+1}}\bigr)
 +\sum_{j=1}^{4}\int m_{a_j}\log m_{a_j}
 -\log a_5 . \qedhere
\]
\end{proof}

\subsection{Summation}

Take the three folds of Section~\ref{sec:finite},
\[
 \phi_1=(0,1,0,1,2,3),\qquad
 \phi_2=(0,1,2,1,2,3),\qquad
 \phi_3=(0,1,2,3,2,3),
\]
and add the three instances of Lemma~\ref{lem:perfold}.  The multiplicities are
those already counted there: each edge of $P_3$ is traversed $3+1+1=5$ times in
total, and among the twelve internal vertices the labels $0,1,2,3$ occur
$1,5,5,1$ times.  Hence
\[
 3\log a_5
 \ \geq\ 5\bigl(E_{01}+E_{12}+E_{23}\bigr)+V_0+5V_1+5V_2+V_3
 \ =\ 5\log a_3+V_0+V_3
\]
by Lemma~\ref{lem:identity}, and $V_0,V_3\geq0$ by
\eqref{eq:rel-entropy-nonneg}.  Therefore
\begin{equation}
 a_5^3\ \geq\ a_3^5
 \qquad\text{whenever }W\geq\varepsilon>0 .
 \label{eq:regularized}
\end{equation}

\subsection{Removing the regularization}
\label{sec:limit}

For a general graphon $W$ and $0<\varepsilon<1$ put
$W_\varepsilon=(1-\varepsilon)W+\varepsilon$, again a graphon, with
$W_\varepsilon\geq\varepsilon$.  Since
$\lVert W_\varepsilon-W\rVert_\infty\leq2\varepsilon$ and a homomorphism
density with $m$ edges is $m$-Lipschitz in the supremum norm,
\[
 |t(P_k,W_\varepsilon)-t(P_k,W)|\leq2k\varepsilon .
\]
Applying \eqref{eq:regularized} to $W_\varepsilon$ and letting
$\varepsilon\downarrow0$ gives $t(P_5,W)^3\geq t(P_3,W)^5$ for every graphon
$W$ on every probability space.  \qed

\begin{remark}
The two isolated vertices of $F=P_0\sqcup P_0\sqcup P_5\sqcup P_5\sqcup P_5$ do
not appear above.  In Section~\ref{sec:finite} they contribute the factor $n^2$
that compensates $-H(X_0)-H(X_3)$; here $\mu$ is already a probability measure,
$t(P_0,W)=1$, and the same compensation is supplied by $V_0,V_3\geq0$.  This is
the only structural difference between the two proofs.
\end{remark}

\begin{remark}
Every step above is finite-dimensional integration against explicit densities
built from $W$, $d$, $A$, $B$ and $a_3$, together with one scalar inequality
(Lemma~\ref{lem:gibbs}) and one limit.  In the Lean development
Lemma~\ref{lem:gibbs} is \texttt{integral\_mul\_log\_div\_nonpos} of
\texttt{Methods/PureChordal/Gibbs.lean}, the limit of
Section~\ref{sec:limit} uses \texttt{Methods/PureChordal/Regularization.lean},
and the coordinate integrations of Lemma~\ref{lem:fold} use the
assignment-integral peeling ladder.  The same combination already carries the
entropy proof of the pure-chordal theorem in
\texttt{Methods/PureChordal/EntropyGluing.lean}.
\end{remark}

%%%%%%%%%%%%%%%%%%%%%%%%%%%%%%%%%%%%%%%%%%%%%%%%%%%%%%%%%%%%%%%%%%%%%%%%%%%%%%%
\appendix
\section{What the Lean formalization actually proves}
%%%%%%%%%%%%%%%%%%%%%%%%%%%%%%%%%%%%%%%%%%%%%%%%%%%%%%%%%%%%%%%%%%%%%%%%%%%%%%%

Section~\ref{sec:graphon} is formalized in full, as
\texttt{a3\_pow\_five\_le\_a5\_pow\_three} in
\texttt{lean/Taeyoung/Methods/OddWalk/}, axiom-clean.  One difference of
substance.

\subsection*{No measure on $\Omega^4$ or $\Omega^6$ is constructed}

The argument above is written against the probability density
\eqref{eq:tree-factorization} on $\Omega^4$, and Lemma~\ref{lem:fold}
integrates a six-coordinate quantity.  Neither space is built.  Everything is
phrased instead inside $\Omega$, through the iterate
$T_W^{\,n}\mathbf1$, with $a_3=\int B$ and $a_5=\int AB$ by self-adjointness of
$T_W$.  Three consequences:

\begin{itemize}
\item The marginals \eqref{eq:marginals} are supplied as four closed-form
 functions --- $P_3$ has two kinds of vertex and two kinds of edge --- rather
 than as marginals of anything.  The remaining two are these with their
 arguments interchanged.
\item The step inequality of Lemma~\ref{lem:perfold} is applied
 \emph{pointwise in the outer variable}: for fixed $y$ the conditional density
 $z\mapsto K(y,z)/u(y)$ is a probability density, which is the Markov property
 of the folded walk, so each step is one application of Lemma~\ref{lem:gibbs}
 on $\Omega$ and one Fubini on $\Omega^2$.  This is what replaces the six-fold
 construction that Section~\ref{sec:finite} performs when it samples a random
 homomorphism.
\item The tree-entropy identity of Lemma~\ref{lem:identity} is not obtained by
 evaluating $\int f\log f$ in two ways.  Because the marginals are explicit
 quotients, every logarithm splits and the identity becomes scalar algebra in
 \[
  P_d=\int_\Omega\frac{dA}{a_3}\log d,\qquad
  P_A=\int_\Omega\frac{dA}{a_3}\log A :
 \]
 the $\log W$ cancels inside each edge term, leaving
 $E_{01}=E_{23}=\log a_3-P_A$, $E_{12}=\log a_3-2P_d$ and
 $V_1=V_2=P_d+P_A-\log a_3$, whose combination is $\log a_3$.
\end{itemize}

Accordingly the closing remark of Section~\ref{sec:limit} is corrected on one
point: the coordinate integrations of Lemma~\ref{lem:fold} do \emph{not} use
the assignment-integral peeling ladder, which the development never invokes.

\end{document}
